%% LyX 2.3.6.1 created this file.  For more info, see http://www.lyx.org/.
%% Do not edit unless you really know what you are doing.
\documentclass[11pt,american,czech]{book}
\usepackage[T1]{fontenc}
\usepackage[utf8]{inputenc}
\usepackage[a4paper]{geometry}
\geometry{verbose,tmargin=4cm,bmargin=3cm,lmargin=3cm,rmargin=2cm,headheight=0.8cm,headsep=1cm,footskip=0.5cm}
\pagestyle{headings}
\setcounter{secnumdepth}{3}
\usepackage{url}
\usepackage{amsmath}
\usepackage{amsthm}
\usepackage{amssymb}
\usepackage{graphicx}
\usepackage{setspace}

\makeatletter
%%%%%%%%%%%%%%%%%%%%%%%%%%%%%% Textclass specific LaTeX commands.
\newenvironment{lyxlist}[1]
	{\begin{list}{}
		{\settowidth{\labelwidth}{#1}
		 \setlength{\leftmargin}{\labelwidth}
		 \addtolength{\leftmargin}{\labelsep}
		 \renewcommand{\makelabel}[1]{##1\hfil}}}
	{\end{list}}

%%%%%%%%%%%%%%%%%%%%%%%%%%%%%% User specified LaTeX commands.
%% Font setup: please leave the LyX font settings all set to 'default'
%% if you want to use any of these packages:

%% Use Times New Roman font for text and Belleek font for math
%% Please make sure that the 'esint' package is turned off in the
%% 'Math options' page.
\usepackage[varg]{txfonts}

%% Use Utopia text with Fourier-GUTenberg math
%\usepackage{fourier}

%% Bitstream Charter text with Math Design math
%\usepackage[charter]{mathdesign}

%%---------------------------------------------------------------------

%% Make the multiline figure/table captions indent so that the second
%% line "hangs" right below the first one.
%\usepackage[format=hang]{caption}

%% Indent even the first paragraph in each section
\usepackage{indentfirst}

%%---------------------------------------------------------------------

%% Disable page numbers in the TOC. LOF, LOT (TOC automatically
%% adds \thispagestyle{chapter} if not overriden
%\addtocontents{toc}{\protect\thispagestyle{empty}}
%\addtocontents{lof}{\protect\thispagestyle{empty}}
%\addtocontents{lot}{\protect\thispagestyle{empty}}

%% Shifts the top line of the TOC (not the title) 1cm upwards 
%% so that the whole TOC fits on 1 page. Additional page size
%% adjustment is performed at the point where the TOC
%% is inserted.
%\addtocontents{toc}{\protect\vspace{-1cm}}

%%---------------------------------------------------------------------

% completely avoid orphans (first lines of a new paragraph on the bottom of a page)
\clubpenalty=9500

% completely avoid widows (last lines of paragraph on a new page)
\widowpenalty=9500

% disable hyphenation of acronyms
\hyphenation{CDFA HARDI HiPPIES IKEM InterTrack MEGIDDO MIMD MPFA DICOM ASCLEPIOS MedInria}

%%---------------------------------------------------------------------

%% Print out all vectors in bold type instead of printing an arrow above them
\renewcommand{\vec}[1]{\boldsymbol{#1}}

% Replace standard \cite by the parenthetical variant \citep
%\renewcommand{\cite}{\citep}

\makeatother

\usepackage{babel}
\begin{document}
\def\documentdate{7. \v{c}ervence 2021}

%%\def\documentdate{\today}

\pagestyle{empty}
{\centering

\noindent %
\begin{minipage}[c]{3cm}%
\noindent \begin{center}
\includegraphics[width=3cm,height=3cm,keepaspectratio]{Images/TITLE/cvut}
\par\end{center}%
\end{minipage}%
\begin{minipage}[c]{0.6\linewidth}%
\begin{center}
\textsc{\large{}České vysoké učení technické v Praze}{\large{}}\\
{\large{}Fakulta jaderná a fyzikálně inženýrská}
\par\end{center}%
\end{minipage}%
\begin{minipage}[c]{3cm}%
\noindent \begin{center}
\includegraphics[width=3cm,height=3cm,keepaspectratio]{Images/TITLE/fjfi}
\par\end{center}%
\end{minipage}

\vspace{3cm}

\textbf{\huge{}Název práce česky název práce česky název práce česky
název práce česky}{\huge\par}

\vspace{1cm}

\selectlanguage{american}%
\textbf{\huge{}Title of the Work in English Title of the Work in English
Title of the Work in English}{\huge\par}

\selectlanguage{czech}%
\vspace{2cm}

{\large{}Bakalářská práce}{\large\par}

}

\vfill{}

\begin{lyxlist}{MMMMMMMMM}
\begin{singlespace}
\item [{Autor:}] \textbf{Jméno Autora}
\item [{Vedoucí~práce:}] \textbf{prof. Ing. Jméno Školitele, DrSc.}
\item [{Konzultant:}] \textbf{doc. RNDr. Jméno Konzultanta, CSc. }(pouze
pokud konzultant byl jmenován.)
\item [{Akademický~rok:}] 2020/2021
\end{singlespace}
\end{lyxlist}
\newpage{}

~\newpage{}

~

\vfill{}

\begin{center}
- Zadání práce -
\par\end{center}

\vfill{}

~\newpage{}

~

\vfill{}

\begin{center}
- Zadání práce (zadní strana) -
\par\end{center}

\vfill{}

~\newpage{}

\noindent \emph{\Large{}Poděkování:}{\Large\par}

\noindent Chtěl bych zde poděkovat především svému školiteli ...................
za pečlivost, ochotu, vstřícnost a odborné i lidské zázemí při vedení
mé diplomové práce. Dále děkuji svému konzultantovi ................
za ................

\vfill

\noindent \emph{\Large{}Čestné prohlášení:}{\Large\par}

\noindent Prohlašuji, že jsem tuto práci vypracoval samostatně a uvedl
jsem všechnu použitou literaturu.

\bigskip{}

\noindent V Praze dne \documentdate\hfill{}Jméno Autora

\vspace{2cm}

\newpage{}

~\newpage{}

\begin{onehalfspace}
\noindent \emph{Název práce:}

\noindent \textbf{Název práce}
\end{onehalfspace}

\bigskip{}

\noindent \emph{Autor:} Jméno Autora

\bigskip{}

\noindent \emph{Obor:} Celý název oboru (nikoliv zkratka)\bigskip{}

\noindent \emph{Zaměření:} Celý název zaměření (Pokud obor neobsahuje
zaměření, tuto řádku odstranit.)

\bigskip{}

\noindent \emph{Druh práce:} Bakalářská práce

\bigskip{}

\noindent \emph{Vedoucí práce:} prof. Ing. Jméno Školitele, DrSc.,
pracoviště školitele (název instituce, fakulty, katedry...)

\bigskip{}

\noindent \emph{Konzultant:} doc. RNDr. Jméno Konzultanta, CSc., pracoviště
konzultanta. Pouze pokud konzultant byl jmenován.

\bigskip{}

\noindent \emph{Abstrakt:} Abstrakt max. na 10 řádků. Abstrakt max.
na 10 řádků. Abstrakt max. na 10 řádků. Abstrakt max. na 10 řádků.
Abstrakt max. na 10 řádků. Abstrakt max. na 10 řádků. Abstrakt max.
na 10 řádků. Abstrakt max. na 10 řádků. Abstrakt max. na 10 řádků.
Abstrakt max. na 10 řádků. Abstrakt max. na 10 řádků. Abstrakt max.
na 10 řádků. Abstrakt max. na 10 řádků. Abstrakt max. na 10 řádků.
Abstrakt max. na 10 řádků. Abstrakt max. na 10 řádků. Abstrakt max.
na 10 řádků. Abstrakt max. na 10 řádků. Abstrakt max. na 10 řádků.
Abstrakt max. na 10 řádků. Abstrakt max. na 10 řádků. Abstrakt max.
na 10 řádků. Abstrakt max. na 10 řádků. Abstrakt max. na 10 řádků.
Abstrakt max. na 10 řádků. Abstrakt max. na 10 řádků. Abstrakt max.
na 10 řádků. Abstrakt max. na 10 řádků. Abstrakt max. na 10 řádků. 

\bigskip{}

\noindent \emph{Klíčová slova:} klíčová slova (nebo výrazy) seřazená
podle abecedy a oddělená čárkou

\vfill{}
~

\selectlanguage{american}%
\begin{onehalfspace}
\noindent \emph{Title:}

\noindent \textbf{Title of the Work}
\end{onehalfspace}

\bigskip{}

\noindent \emph{Author:} Author's Name

\bigskip{}

\noindent \emph{Abstract:} Max. 10 lines of English abstract text.
Max. 10 lines of English abstract text. Max. 10 lines of English abstract
text. Max. 10 lines of English abstract text. Max. 10 lines of English
abstract text. Max. 10 lines of English abstract text. Max. 10 lines
of English abstract text. Max. 10 lines of English abstract text.
Max. 10 lines of English abstract text. Max. 10 lines of English abstract
text. Max. 10 lines of English abstract text. Max. 10 lines of English
abstract text. Max. 10 lines of English abstract text. Max. 10 lines
of English abstract text. Max. 10 lines of English abstract text.
Max. 10 lines of English abstract text. Max. 10 lines of English abstract
text. Max. 10 lines of English abstract text. Max. 10 lines of English
abstract text. Max. 10 lines of English abstract text. Max. 10 lines
of English abstract text. Max. 10 lines of English abstract text.
Max. 10 lines of English abstract text. Max. 10 lines of English abstract
text. Max. 10 lines of English abstract text.

\bigskip{}

\noindent \emph{Key words:} keywords in alphabetical order separated
by commas

\selectlanguage{czech}%
\newpage{}

~\newpage{}

\pagestyle{plain}

\tableofcontents{}

\newpage{}

\chapter*{Úvod}

\addcontentsline{toc}{chapter}{Úvod}

Text úvodu....

\chapter{Název první kapitoly}

\pagestyle{headings}

\chapter*{Závěr}

\pagestyle{plain}

\addcontentsline{toc}{chapter}{Záv\v{e}r}

Text závěru....
\begin{thebibliography}{1}
\bibitem{Allen-Cahn}S. Allen, J. W. Cahn: \emph{A microscopic theory
for antiphase boundary motion and its application to antiphase domain
coarsening}. Acta Metall., 27:1084-1095, 1979.

\bibitem{CINECA}G. Ballabio et al.: \emph{High Performance Systems
User Guide}. High Performance Systems Department, CINECA, Bologna,
2005. \url{www.cineca.it}

\bibitem{rumpf3}J. Becker, T. Preusser, M. Rumpf: \emph{PDE methods
in flow simulation post processing}. Computing and Visualization in
Science, 3(3):159-167, 2000.
\end{thebibliography}

\end{document}
